\scanpage{150}
\mspage{76}
% archive sheet number: 149
% L1
\deleted{est non triviale. Donc pour voir si}

% L2–5
Notons que l'on obtient déjà la valeur ``universelle'' de $(\mathrm{Gl}'_{\deleted{\unsure{n}}}(2,\mathbb Z/n\mathbb Z)^{\pm})_{\rm ab}$, \inserted{savoir $\mathbb F_2\times\mathbb F_2$,} pour $n=4$, de telle façon qu'il est quand $=$ évident a priori que $\varepsilon$ se factorise par
% L6
\[
\widehat{\mathbb Z}^{*}/\pm1\mathrel{\unsure{-|\longrightarrow}}
\underbrace{(\deleted{\unsure{\widehat Z}}\mathbb Z/\deleted{\unsure{8}}4\mathbb Z)^{*}/\pm1}_{\substack{\deleted{\unsure{\pm1\simeq Z'_8}}\\Z'_4=\{1\}}}
\longrightarrow\{\pm1\}\,,
\tag*{\deleted{(\unsure{55|56})}}
\]
% L7–11
\deleted{et la seule question est si ce caractère est ou non trivial, i.e. si $(\mathbb Z/8\mathbb Z)^{*}/\pm1$ \inserted{$\simeq\pm1$} opère trivialement, ou non, sur $(\mathrm{Gl}'(2,\mathbb Z/8\mathbb Z)^{\pm})_{\rm ab}\simeq\mathbb F_2\times\mathbb F_2$.}

\deleted{Notons que $\mathrm{Gl}'(2,\mathbb Z/4\mathbb Z)^{\pm}\simeq\unsure{\mathfrak S_4|\mathscr S_4}\times\pm1$.}
% L12
ce qui montre que \uline{le caractère $\varepsilon$ est trivial !}$]\!]$

% L13–17
\deleted{\unsure{Reven|Revenons}} Notons que \inserted{\unsure{si $n$|si on}} $G_0^+=\mathrm{Aut}^+(S,F_{i_0})$ est le seul sous-groupe d'indice 2 de $G_0$ qui contienne $\rho_0$ (\uline{NB} $\dot\rho_0$ \deleted{\unsure{$=\dot r_0\dot r_\infty$}} $\dot r_1\dot r_\infty$) \deleted{\unsure{et}}\inserted{-- où $\rho_0$ est l'opération de rotation autour d'un sommet déterminé $\deleted{\unsure{a}}s_0\in S$, associée au repère \unsure{$(a_0,s_0,s'_0)$|$(s_0,s_1,s_\infty)$}}, pour prouver que \deleted{\unsure{il est}} $G_0^+$ égal à $G_1^+=\mathrm{Aut}^+(S,F_{i_1})$, pour $i_1=\lambda.i_0$, il suffit de voir que l'on a
% L18
\[
\rho_0\overset{\mathrm{déf}}{=}\rho_{s_0,A_0}\in G_1^+=\mathrm{Aut}^+(S,F_i)\,,
\tag*{(55)}
\]
% L19–20
pour ceci considérons l'étoile
\[
E=\mathrm{Et}_0(s_0)=\mathrm{Et}_1(\dot\lambda.s_0)
\]
% L21–25
$\rho_0$ applique $E$ dans lui-$=$, et est un automorphisme direct \inserted{i.e. une ``rotation''} pour la structure polygonale $A_{0E}=A_0\cap\unsure{\mathcal P_2|P_2}(E)$, \deleted{\unsure{On en devrait aussitôt}}\deleted{\unsure{i.e. appartient à}} On en conclut qu'il est aussi direct pour la structure polygonale $A_{1E}=A_1\cap\mathcal P_2(E)$, puisque $A_{1E}=A_{0E}^{(\dot\lambda)}$ -- donc
